\chapter{Introdução}
\label{cap:introducao}

	Para onde vamos? uma das questões fundamentais da natureza humana, estudada pelas mais diferentes ciências, desde exatas às sociais. Dentre as ciências que buscam a resposta do para onde iremos, temos a economia, uma ciência que encontra-se no meio termo, a chamada social aplicada. Inicialmente pode parecer estranho, mas boa parte do trabalho de um economista é entender os fatos passados e presentes para predizer (ou tentar) o  futuro, mas a pergunta que fica é por quê a economia se preocupa tanto com o futuro?

Para \citeonline[p. 1]{mankiw2025introduccao} "a economia é o estudo de como a sociedade administra seus recursos escassos", é essa escassez \footnote{escassez:  natureza limitada dos recursos da sociedade \citem[p. 1]{mankiw2025introduccao}} que é o motivador para a preocupação com o futuro, afinal se tivéssemos recursos ilimitados não haveria motivo para estudar economia. Mas este não é o caso, existe limitação de recursos, e devemos saber qual é a melhor maneira de alocar esses recursos.

Para um melhor enfoque, a economia se subdivide em duas grandes áreas, a microeconomia, voltada para o estudo da firma (empresas), e a macroeconomia, que estuda os agregados econômicos \citem{Vasconcelos}

No presente trabalho nos voltaremos para a macroeconomia, área da economia que surge no inicio do século XX, devido as intensas mudanças sociais, políticas e econômicas que ocorreram no período, a principal obra referência da macroeconomia é o livro de John Maynard Keynes: \textit{Teoria Geral do Emprego, do Juro e da Moeda} \citeyear{KEYNES1936}, onde o autor traz os principais fundamentos da macroeconomia que é estudada hoje na academia e executada nos governos.


Um desses conceitos fundamentais que Keynes aborda é a noção de incerteza
\begin{citacao}
	De fato, sabemos muito pouco sobre o futuro. A base real para formar expectativas sobre o rendimento futuro de um investimento é extremamente frágil. Nossas decisões de investimento [...] dependem de um estado de confiança. Muitas vezes, sabemos que não sabemos; mas mesmo assim temos que agir. \citem[p. 161]{KEYNES1936}
\end{citacao}
Com essa noção de incerteza definida por Keynes podemos relacionar com aquela nossa questão fundamental, gestores mundiais, sejam eles públicos ou privados, buscam maneiras de observar dados históricos e com um certo intervalo de confiança tentar "prever o futuro", para justamente, combater a incerteza.


A incerteza 
%	traz desconfiança, e a desconfiança, por sua vez, 
gera perguntas, e encontrar as respostas para essas perguntas vai muito além da impressão pessoal
%	\textit{feeling}
%	\footnote{\textit{feeling}: impressão, sentimento, pressentimento}
do gestor. Os gestores precisam de informação para suas decisões, e para \citeonline{laudon2022sistemas} os sistemas de informação são essenciais para alcançar objetivos estratégicos dos negócios. Na era digital, o grande volume de dados gerados pelos sistemas de informação podem ser usados como fontes, praticamente inesgotáveis de dados.

Esses grandes volumes de dados, gerados pelos sistemas de informação, são chamados de \textit{big data}. Para \citem[p. 9]{vasconcelos2017poder}  "O \textit{big data} caracterizado por grandes bases de dados com um número elevado de preditores relativos às observações disponíveis". Contudo, somente ter os dados não é o suficiente, e converter esse grande volume de dados em informações não é uma tarefa simples, e exige técnicas sofisticadas de coleta, tratamento e modelagem. Os processos de coleta e tratamento, apesar de fundamentais, serão usados somente como passos anteriores ao da modelagem, que é o enfoque  desse texto.  

% gostaria de trazer um autor sobre modelos econometricos, porém fiquei indeciso













%	Nesse processo a estatística e a ciência de dados trazem o aporte teórico para o tratamento, coleta e transformação desses dados em informações uteis. 


%	Adicioanar a definição de IA, Machine Learning na Introdução não esquecer


\section{Problema de pesquisa}

Avaliar modelos na previsão de variáveis econômicas não é um processo novo, muitos autores já usaram diversas vezes esses modelos.

Em sua dissertação, \citeonline{azevedo2020modelagem} usou o modelo modelo de vetores auto-regressivos (VAR) para análise de séries temporais com variáveis macroeconômicas.
Já \citeonline{fialho2024incerteza} usaram o VAR para entender a dinâmica entre 
incerteza, spread bancário, crédito, consumo, rendimento e taxa de desocupação. 

\citeonline{ciquini2023modelo} avaliou o uso do modelo auto-regressivo de médias móveis (do inglês \textit{AutoRegressive Integrated Moving Average} - ARIMA) para tomada de decisões econômicas. Agora \citeonline{Silva_Clericuzi_Carvalho_2025} apoiaram-se no ARIMA para a predição do PIB brasileiro. 

\citeonline{alves2009nucleo} em sua dissertação de mestrado abordou modelos de fatores dinâmicos (do inglês  \textit{Dynamic Factor Models} - DFM) para entender o núcleo da inflação brasileira. Por outro lado \citeonline{moura2020fragilidade} avaliou os ciclos econômicos brasileiros usando uma combinação entre o DFM e o VAR.

\citeonline{castelli2022modelo} abordou o uso do modelo de amostragem de dados mistos (do inglês \textit{Mixed Data Sampling} - MIDAS) para predizer o PIB com a pesquisa industrial mensal. Assim como \citeonline{micheletti2024avaliaccao} usou modelos de \textit{ensemble stacking} \footnote{\textit{ensemble stacking}: modelos combinados} (dentre eles o MIDAS) para prever o PIB.

A maioria desses autores trabalharam com o que se chamam modelos tradicionais de predição, esses modelos se baseiam em variáveis e proposições bem definidas, o que pode não ser o caso na era do \textit{big data}. Nesse contexto, em que os modelos tradicionais encontram-se limitados, que a inovação surge, nesse caso o uso de técnicas de IA para modelar variáveis econômicas.

Diante do disso a presente pesquisa tende a seguir os passos de \citeonline{kava2022alem} e de \citeonline{rebelo2022interpretabilidade} que comparam, em suas dissertações, modelos de IA com os chamados modelos tradicionais. Contudo, a nossa pesquisa se diferencia dos trabalhos produzidos por testar outros modelos e focar especificamente no PIB brasileiro, para isso norteia-se pela seguinte pergunta: Em que medida modelos preditivos baseados em \ac{ia}, comparados aos modelos tradicionais VAR, ARIMA, DFM e MIDAS, conseguem estimar o PIB brasileiro?

%Como desenvolver um modelo preditivo baseado em IA para efetivamente estimar i]o PIB brasileiro, considerando as diferenças e semlhanças significativas entre os modelos tradicionais VAR, ARIMA, DFM e MIDAS.

\section{Objetivos}
\subsection{Objetivo Geral}
Analisar em que medida modelos preditivos baseados em \ac{ia}, comparados aos modelos tradicionais VAR, ARIMA, DFM e MIDAS, conseguem estimar o PIB brasileiro.

\subsection{Objetivos Específicos}
\begin{itemize}
	\item[a)] Caracterizar os modelos econométricos tradicionais (VAR, ARIMA, DFM e MIDAS) e de IA para a predição do PIB;
	\item[b)] Identificar as diferentes variáveis determinantes na estimativa do PIB;
	\item[c)] Desenvolver um ciclo de vida para a produção de um novo modelo econométrico de IA para a predição do PIB; 
	\item[d)] Comparar o desempenho dos modelos de IA e econométricos tradicionais usando diferentes métricas, como desempenho computacional, precisão e interpretabilidade.
\end{itemize}



\section{Justificativa} 
O presente trabalho justifica-se empiricamente, pois permite que gestores, sejam eles públicos ou privados, utilizem-se de ferramentas estatísticas para tomarem decisões fundamentadas. "A  tomada  de  decisões  empresariais  é  uma  tarefa  complexa  e  fundamental  para  a  sobrevivência de uma organização" \citem[p. 3]{santana2023aprimorando}.

Portanto, a tomada de decisões pode ser fundamentada através dos modelos. Os modelos, mesmo que sejam uma simplificação da realidade, podem aumentar a confiança dos agentes. Segundo \citeonline{KEYNES1936}, essa confiança é um fator que afeta as decisões de investimento.

Já contribuição teórica desse trabalho está alinhada com a avaliação dos diferentes modelos estatísticos para a predição econômica, além é claro, de explorar as novas fronteiras da ciência, com a avaliação de modelos inovadores como os que usam a IA. 


O presente trabalho busca mesclar diferentes modelos, ranqueando seu desempenho, e possibilitando a identificação dos modelos mais adequados para previsão econômica. Para isso é necessário a avaliação métricas estatísticas comuns, semelhante ao que \citeonline{castelli2022modelo} fez, ao comparar modelos de IA com modelos tradicionais, aplicando métricas como coeficiente de determinação ($R^2$), Erro Médio Absoluto (MAE), Erro Quadrático Médio (MSE) para classificar os modelos.

	
